
\section{Fundamento teórico}
\subsection{Ecuaciones de Maxwell en medios lineales}

El comportamiento de la radiación electromagnética en un medio lineal, isótropo 
y homogéneo está gobernado por las ecuaciones de Maxwell en ausencia de cargas y 
corrientes libres:

% \begin{equation}
%   \nabla \times \mathbf{E} = -\mu \frac{\partial \mathbf{H}}{\partial t}.
% \end{equation}
% \begin{equation}
%   \nabla \times \mathbf{H} =  \epsilon \frac{\partial \mathbf{E}}{\partial t}.
% \end{equation}
% \begin{equation}
%   \nabla \cdot (\epsilon \mathbf{E}) = 0.
% \end{equation}
% \begin{equation}
%   \nabla \cdot (\mu \mathbf{H}) = 0.
% \end{equation}


\begin{equation}
\begin{aligned}
\nabla \times \mathbf{E} &= -\mu \frac{\partial \mathbf{H}}{\partial t}
&\qquad
\nabla \cdot (\epsilon \mathbf{E}) &= 0 \\
\nabla \times \mathbf{H} &=  \epsilon \frac{\partial \mathbf{E}}{\partial t}
&\qquad
\nabla \cdot (\mu \mathbf{H}) &= 0
\end{aligned}
\end{equation}

donde $\mathbf{E}$ y $\mathbf{H}$ son los campos eléctrico y magnético, respectivamente, 
y $\epsilon$ y $\mu$ son la permitividad y permeabilidad del medio. 
% En medios no magnéticos 
% se suele tomar $\mu \approx \mu_0$, mientras que la respuesta de la materia se incorpora 
% en la permitividad efectiva $\epsilon$.

Suponiendo una dependencia temporal armónica de la forma $\exp(-i\omega t)$, las ecuaciones 
de Maxwell se reducen a la ecuación de onda vectorial para el campo eléctrico y magnético:

\begin{equation}
  \nabla^2 \mathbf{E} + k^2 \mathbf{E} = 0, \ \ \ \nabla^2 \mathbf{H} + k^2 \mathbf{H} = 0,
\end{equation}

donde

\begin{equation}
  k = \frac{2\pi}{\lambda} n
\end{equation}

es el número de onda en el medio, $\lambda$ la longitud de onda en el vacío y $n$ el índice 
de refracción del medio circundante. En materiales absorbentes, el índice de refracción es 
complejo, $m = n + i k$, y la parte imaginaria da cuenta de la absorción interna.

El problema de la dispersión de la luz por una partícula esférica consiste en resolver 
la ecuación de onda en dos dominios: el interior de la esfera, caracterizado por un índice 
de refracción $m$, y el exterior, caracterizado por el índice del medio circundante 
(en este trabajo, vacío). Las soluciones deben satisfacer las condiciones de contorno 
en la superficie de la partícula, que exigen la continuidad de las componentes tangenciales 
de los campos eléctrico y magnético.


Gracias a la linealidad de las ecuaciones de Maxwell y de las condiciones de 
contorno, cualquier solución física del problema puede construirse como una 
superposición de soluciones fundamentales. Esto significa que la dispersión 
de un campo incidente arbitrario puede obtenerse combinando adecuadamente 
las soluciones asociadas a ondas planas con polarizaciones ortogonales. 
Esta propiedad justifica que la teoría de Mie resuelva únicamente el problema 
para dos polarizaciones independientes, ya que cualquier polarización del campo 
incidente puede expresarse como una suma lineal de ambas.

\subsection{Condiciones de contorno y matriz de amplitud de dispersión}

El problema de la dispersión consiste en resolver la ecuación de onda en 
dos dominios: el interior de la esfera, caracterizado por un índice de 
refracción $m$, y el exterior, caracterizado por el índice del medio 
circundante. Las soluciones en ambos dominios deben satisfacer las 
condiciones de contorno en la superficie $r = a$, que exigen la 
continuidad de las componentes tangenciales de los campos eléctrico y 
magnético:

\begin{equation}
  \left.(E_\theta^{(i)} + E_\theta^{(s)})\right|_{r=a} 
  = \left.E_\theta^{(\mathrm{int})}\right|_{r=a},
  \qquad
  \left.(H_\theta^{(i)} + H_\theta^{(s)})\right|_{r=a} 
  = \left.H_\theta^{(\mathrm{int})}\right|_{r=a},
\end{equation}

y análogamente para las componentes azimutales $\phi$. 

Dado que tanto las ecuaciones de Maxwell como estas condiciones de contorno 
son lineales en los campos, el campo dispersado en cualquier punto del 
espacio debe depender linealmente del campo incidente. En la zona de campo 
lejano ($kr \gg 1$), el campo dispersado adopta la forma de 
onda esférica saliente Jackson 1975 pag 270:

\begin{equation}
  \mathbf{E}^{(s)} \;\xrightarrow{kr \to \infty}\; 
  \frac{e^{ikr}}{-ikr}\,\mathbf{f}(\theta,\phi),
\end{equation}

donde $\mathbf{f}$ es la amplitud vectorial de dispersión. Combinando la 
linealidad con esta estructura de campo lejano, la relación entre las 
componentes de polarización del campo incidente y del campo dispersado debe 
ser de la forma matricial:

\begin{equation}
\label{eq:amplitudematrix}
  \begin{pmatrix}
    E_{\parallel}^{(s)} \\[4pt]
    E_{\perp}^{(s)}
  \end{pmatrix}
  =
  \frac{e^{ikr}}{-ikr}
  \begin{pmatrix}
    S_{2}(\theta,\phi) & S_{3}(\theta,\phi) \\[4pt]
    S_{4}(\theta,\phi) & S_{1}(\theta,\phi)
  \end{pmatrix}
  \begin{pmatrix}
    E_{\parallel}^{(i)} \\[4pt]
    E_{\perp}^{(i)}
  \end{pmatrix}.
\end{equation}

Aquí, $E_{\parallel}^{(i)}$ y $E_{\perp}^{(i)}$ son las componentes paralela y perpendicular
del campo eléctrico incidente (paralelo y perpendicular respecto al plano de incidencia), 
mientras que $E_{\parallel}^{(s)}$ y $E_{\perp}^{(s)}$ son las
componentes correspondientes del campo dispersado (\textit{scattered}). 
Los elementos $S_1, S_2, S_3, S_4$ dependen del ángulo de dispersión 
$\theta$ y del ángulo azimutal $\phi$, y quedan completamente determinados 
por las condiciones de contorno sobre la superficie de la partícula. Para 
una esfera homogénea e isotrópica, la simetría azimutal impone $S_3 = S_4 = 0$, 
de modo que los modos de polarización paralelo y perpendicular se dispersan 
de forma independiente.


============================================


La formulación estándar introduce así la matriz de amplitud de dispersión 
(\ref{eq:amplitudematrix}), que relaciona de forma lineal los campos incidente 
y dispersado y permite caracterizar completamente la interacción entre 
la onda y la partícula. Esta matriz surge directamente de aplicar las condiciones de continuidad 
del campo eléctrico y magnético sobre la superficie esférica, y constituye 
la base para calcular magnitudes observables como los patrones angulares de 
dispersión, la sección eficaz o el grado de polarización del campo dispersado.

\begin{equation}
\label{eq:amplitudematrix}
  \begin{pmatrix}
E_{\parallel}^{(s)} \\
E_{\perp}^{(s)}
\end{pmatrix}
=
\frac{e^{ikr}}{-ikr}
\begin{pmatrix}
S_{2}(\theta, \phi) & S_{3}(\theta, \phi) \\
S_{4}(\theta, \phi) & S_{1}(\theta, \phi)
\end{pmatrix}
\begin{pmatrix}
E_{\parallel}^{(i)} \\
E_{\perp}^{(i)}
\end{pmatrix}.
\end{equation}

Aquí, $E_{\parallel}^{(i)}$ y $E_{\perp}^{(i)}$ son las componentes paralela y perpendicular
del campo eléctrico incidente, mientras que $E_{\parallel}^{(s)}$ y $E_{\perp}^{(s)}$ son las
componentes correspondientes del campo dispersado (\textit{scattered}). Los elementos de la matriz de amplitud
de dispersión, $S_{1}(\theta, \phi)$, $S_{2}(\theta, \phi)$, $S_{3}(\theta, \phi)$ y 
$S_{4}(\theta, \phi)$, dependen del ángulo de dispersión $\theta$ y angulo acimutal $\phi$, y
contienen toda la información sobre cómo la partícula afecta a la onda incidente.


%====================================================================%
%====================================================================%

% \subsection{Expansión en armónicos esféricos vectoriales}
% \hspace{1em}La ecuación de onda vectorial para el campo eléctrico, junto con la condición
% de transversidad $\nabla \cdot \mathbf{E} = 0$, describe modos 
% electromagnéticos físicamente realizables en el medio. En el caso de una 
% partícula esférica, el sistema posee simetría esférica, por lo que el 
% sistema de coordenadas natural es el esférico $(r,\theta,\phi)$. En este 
% sistema, las soluciones separables de la ecuación de onda escalar
% \begin{equation}
%   \nabla^2 \psi + k^2 \psi = 0
% \end{equation}
% pueden escribirse como
% \begin{equation}
%   \psi(r,\theta,\phi) = R(r)\,\Theta(\theta)\,\Phi(\phi),
% \end{equation}
% donde $R(r)$ está dado por funciones de Bessel esféricas y la dependencia
% angular viene dada por polinomios de Legendre asociados.

% Para construir campos electromagnéticos vectoriales que satisfagan tanto 
% la ecuación de onda como la condición de divergencia nula, se introduce 
% la siguiente construcción a partir de una función escalar $\psi$ y un 
% vector constante $\mathbf{c}$:
% \begin{equation}
%   \mathbf{M} = \nabla \times (\psi\,\mathbf{c}), \qquad
%   \mathbf{N} = \frac{1}{k}\,\nabla \times \mathbf{M}.
% \end{equation}

% Los campos $\mathbf{M}$ y $\mathbf{N}$ tienen todas las propiedades necesarias
% para representar campos electromagnéticos físicos: Satisfacen la ecuación de onda
% vectorial, la condición de divergencia nula, el rotacional de $\mathbf{M}$ es 
% proporcional a $\mathbf{N}$ y viceversa. Por lo tanto, el problema para encontrar 
% soluciones para las ecuaciones de campo se reduce a encontrar soluciones a la función
% de onda escalar $\psi$.

% Los campos $\mathbf{M}$ y $\mathbf{N}$ 
% %satisfacen la ecuación de onda vectorial y son ortogonales a 
% %$\nabla \psi$, de modo que 
% describen modos transversos adecuados para representar campos eléctricos y magnéticos.
% A partir de las soluciones escalares separables, esta construcción conduce 
% a las familias de armónicos esféricos vectoriales
% \begin{equation}
%   \mathbf{M}_{nm}(kr,\theta,\phi), \qquad
%   \mathbf{N}_{nm}(kr,\theta,\phi),
% \end{equation}
% que forman un conjunto completo y ortogonal. Cualquier campo 
% electromagnético divergencia–libre en una región con simetría esférica 
% puede expandirse como una superposición de estos modos. En el contexto 
% de la dispersión por una esfera, esta propiedad es esencial, ya que 
% permite escribir de forma consistente los campos incidente, interno 
% y dispersado en la misma base.





% En forma de componentes, los armónicos esféricos vectoriales se expresan como

% \begin{equation}
% \mathbf{M}_{emn}
% = -\frac{m}{\sin\theta}\sin m\phi\, P_n^m(\cos\theta)\, z_n(\rho)\,\hat{\boldsymbol{e}}_\theta
%    - \cos m\phi\, \frac{d P_n^m(\cos\theta)}{d\theta}\, z_n(\rho)\,\hat{\boldsymbol{e}}_\phi .
% \end{equation}

% \begin{equation}
% \mathbf{M}_{omn}
% =  \frac{m}{\sin\theta}\cos m\phi\, P_n^m(\cos\theta)\, z_n(\rho)\,\hat{\boldsymbol{e}}_\theta
%    - \sin m\phi\, \frac{d P_n^m(\cos\theta)}{d\theta}\, z_n(\rho)\,\hat{\boldsymbol{e}}_\phi .
% \end{equation}

% % \begin{equation}
% % \mathbf{N}_{emn}
% % = \frac{z_n(\rho)}{\rho}\cos m\phi\, n(n+1) P_n^m(\cos\theta)\,\hat{\boldsymbol{e}}_r
% % + \cos m\phi\, \frac{d P_n^m(\cos\theta)}{d\theta}\, \frac{1}{\rho}\frac{d}{d\rho}\!\left[\rho z_n(\rho)\right]\hat{\boldsymbol{e}}_\theta
% % - m\sin m\phi\, \frac{P_n^m(\cos\theta)}{\sin\theta}\, \frac{1}{\rho}\frac{d}{d\rho}\!\left[\rho z_n(\rho)\right]\hat{\boldsymbol{e}}_\phi .
% % \end{equation}

% \begin{equation}
% \begin{split}
% \mathbf{N}_{omn} =
% & \frac{z_n(\rho)}{\rho}\sin m\phi\, n(n+1) P_n^m(\cos\theta)\,\hat{\mathbf{e}}_r \\
% & + \sin m\phi\, \frac{d P_n^m(\cos\theta)}{d\theta}\,
%     \frac{1}{\rho}\frac{d}{d\rho}\!\left[\rho z_n(\rho)\right]\hat{\mathbf{e}}_\theta \\
% & + m\cos m\phi\, \frac{P_n^m(\cos\theta)}{\sin\theta}\,
%     \frac{1}{\rho}\frac{d}{d\rho}\!\left[\rho z_n(\rho)\right]\hat{\mathbf{e}}_\phi .
% \end{split}
% \end{equation}



% % \begin{equation}
% % \mathbf{N}_{omn}
% % = \frac{z_n(\rho)}{\rho}\sin m\phi\, n(n+1) P_n^m(\cos\theta)\,\hat{\boldsymbol{e}}_r
% % + \sin m\phi\, \frac{d P_n^m(\cos\theta)}{d\theta}\, \frac{1}{\rho}\frac{d}{d\rho}\!\left[\rho z_n(\rho)\right]\hat{\boldsymbol{e}}_\theta
% % + m\cos m\phi\, \frac{P_n^m(\cos\theta)}{\sin\theta}\, \frac{1}{\rho}\frac{d}{d\rho}\!\left[\rho z_n(\rho)\right]\hat{\boldsymbol{e}}_\phi .
% % \end{equation}


% \begin{equation}
% \begin{split}
% \mathbf{N}_{omn} =
% & \frac{z_n(\rho)}{\rho}\sin m\phi\, n(n+1) P_n^m(\cos\theta)\,\hat{\mathbf{e}}_r \\
% & + \sin m\phi\, \frac{d P_n^m(\cos\theta)}{d\theta}\,
%     \frac{1}{\rho}\frac{d}{d\rho}\!\left[\rho z_n(\rho)\right]\hat{\mathbf{e}}_\theta \\
% & + m\cos m\phi\, \frac{P_n^m(\cos\theta)}{\sin\theta}\,
%     \frac{1}{\rho}\frac{d}{d\rho}\!\left[\rho z_n(\rho)\right]\hat{\mathbf{e}}_\phi .
% \end{split}
% \end{equation}


% En particular, un campo incidente arbitrario, como una onda plana, 
% puede expandirse como
% \begin{equation}
%   \mathbf{E}^{(i)} =
%   \sum_{n=1}^{\infty}\sum_{m=-n}^{n}
%   \left[
%     A_{nm}\,\mathbf{M}_{nm}(kr,\theta,\phi) +
%     B_{nm}\,\mathbf{N}_{nm}(kr,\theta,\phi)
%   \right],
% \end{equation}
% donde los coeficientes $A_{nm}$ y $B_{nm}$ se obtienen utilizando las
% relaciones de ortogonalidad de los armónicos esféricos vectoriales.
% Los campos interno y dispersado admiten expansiones análogas, con 
% coeficientes distintos.

% El uso de esta base tiene dos ventajas fundamentales. En primer lugar, 
% cada término de la expansión satisface de forma automática la ecuación 
% de onda en el dominio correspondiente (interior o exterior a la esfera). 
% En segundo lugar, las condiciones de contorno en la superficie 
% $r = a$ desacoplan los modos: para cada par $(n,m)$ se obtienen 
% ecuaciones algebraicas independientes para los coeficientes de las 
% expansiones interna y externa. Esto hace posible determinar de manera 
% sistemática los coeficientes de dispersión de Mie y, a partir de ellos, 
% calcular las secciones eficaces y los elementos de la matriz de 
% dispersión introducida anteriormente.



%====================================================================%
%====================================================================%

\subsection{Expansión en armónicos esféricos vectoriales}
% Esto es una posible reformulación alternativa de la sección anterior.
\hspace{1em}La resolución del problema de dispersión por una partícula esférica requiere 
expresar los campos electromagnéticos en una base que respete la simetría del 
sistema. La ecuación de onda vectorial
\[
\nabla^2 \mathbf{E} + k^2 \mathbf{E} = 0, \ \ \ \nabla^2 
\mathbf{H} + k^2 \mathbf{H} = 0,\]
junto con la condición de transversidad $\nabla \cdot \mathbf{E} = 0$, 
$\nabla \cdot \mathbf{H} = 0$, describe 
modos electromagnéticos físicamente realizables. Cuando el dominio posee 
simetría esférica, como ocurre en la teoría de Mie, resulta natural buscar 
soluciones que se separen en coordenadas esféricas y que puedan combinarse para 
reproducir los campos incidente, interno y dispersado. La motivación para 
introducir los armónicos esféricos vectoriales surge precisamente de esta 
necesidad: disponer de una familia de funciones base que satisfagan la ecuación 
de onda y se adapten a la geometría del problema.

Un punto de partida consiste en introducir los siguientes campos vectoriales a 
partir de una función escalar arbitraria $\psi$ y de un vector constante 
$\mathbf{r}$:
\[
\mathbf{M} = \nabla \times (\psi\,\mathbf{r}), 
\qquad
\mathbf{N} = \frac{1}{k}\,\nabla \times \mathbf{M},
\]
Utilizando las identidades vectoriales, se consigue
\[
\nabla^2 \mathbf{M} + k^2 \mathbf{M} = \nabla \times [\mathbf{r} 
(\nabla^2 \psi + k^2 \psi)] = 0.\]

Por lo tanto, si la función escalar $\psi$ es una solucion a la ecuación de onda
escalar $\nabla^2 \psi + k^2 \psi = 0$, entonces $\mathbf{M}$ y $\mathbf{N}$
satisfacen la ecuación de onda vectorial. Por lo tanto, ambas estructuras
tienen las propiedades necesarias para describir ondas electromagnéticas: 
ambos vectores satisfacen la ecuación de onda vectorial, poseen divergencia nula 
el rotacional de $\mathbf{M}$ es proporcional a $\mathbf{N}$ y viceversa. 
De este modo, el problema electromagnético queda reducido a encontrar soluciones 
escalares $\psi$ que puedan emplearse como generadoras de campos físicos 
mediante las estructuras $\mathbf{M}$ y $\mathbf{N}$.

La ecuación de onda escalar anterior admite soluciones separables en coordenadas 
esféricas de la forma
\[
\psi(r,\theta,\phi) = R(r)\,\Theta(\theta)\,\Phi(\phi),
\]
lo que conduce a dos ecuaciones angulares y una ecuación radial:
\begin{equation}
\frac{d^{2}\Phi}{d\phi^{2}} + m^{2}\Phi = 0 ,
\end{equation}

\begin{equation}
\frac{1}{\sin\theta}\frac{d}{d\theta}\!\left( \sin\theta \frac{d\Theta}{d\theta} \right)
+ \left[ n(n+1) - \frac{m^{2}}{\sin^{2}\theta} \right] \Theta = 0 ,
\end{equation}

\begin{equation}
\frac{d}{dr}\!\left( r^{2}\frac{dR}{dr} \right)
+ \left[ k^{2}r^{2} - n(n+1) \right] R = 0 .
\end{equation}
La parte 
radial da lugar a funciones de Bessel esféricas, mientras que la dependencia 
angular queda descrita por polinomios de Legendre asociados y funciones 
trigonométricas. Al imponer condiciones de periodicidad en la coordenada 
azimutal y regularidad en el eje polar, aparecen las funciones 

\begin{equation}
\psi_{emn} = \cos(m\phi)\, P_{n}^{m}(\cos\theta)\, z_{n}(kr),
\end{equation}

\begin{equation}
\psi_{omn} = \sin(m\phi)\, P_{n}^{m}(\cos\theta)\, z_{n}(kr).
\end{equation}

donde $P_{n}^{m}$ son los polinomios de Legendre asociados y $z_{n}$ representa 
cualquiera de las funciones de Bessel esféricas $j_n$, $y_n$, $h_n^{(1)}$ o 
$h_n^{(2)}$.

Estas funciones escalares permiten generar los armónicos esféricos vectoriales 
mediante la construcción anterior. Los campos resultantes se expresan como
\begin{equation} 
\label{Generadores vectoriales}
\begin{aligned}
\mathbf{M}_{emn} &= \nabla \times (\psi_{emn}\,\mathbf{r}), 
\qquad
\mathbf{M}_{omn} &= \nabla \times (\psi_{omn}\,\mathbf{r}), 
\\
\mathbf{N}_{emn} &= \frac{1}{k}\,\nabla \times \mathbf{M}_{emn},
\qquad
\mathbf{N}_{omn} &= \frac{1}{k}\,\nabla \times \mathbf{M}_{omn}.
\end{aligned}
\end{equation}

% \textit{y representan modos transversos que forman una base completa y ortogonal en 
% regiones con simetría esférica. Esta propiedad es esencial en la teoría de Mie: 
% al expandir el campo incidente, el interno y el dispersado en esta misma base, 
% se consigue que las condiciones de contorno sobre la superficie $r = a$ 
% desacoplen los modos para cada par $(n,m)$, lo que permite determinar de forma 
% directa los coeficientes de Mie.} no me gusta esta frase, pero la dejo por si acaso.

% En forma explícita, los armónicos esféricos vectoriales pueden escribirse en 
% componentes angulares y radiales como


% \begin{equation}\label{eq:Memn}
% \mathbf{M}_{emn}
% = -\frac{m}{\sin\theta}\sin m\phi\, P_n^m(\cos\theta)\, z_n(\rho)\,\hat{\boldsymbol{e}}_\theta
%   - \cos m\phi\, \frac{d P_n^m(\cos\theta)}{d\theta}\, z_n(\rho)\,\hat{\boldsymbol{e}}_\phi ,
% \end{equation}

% \begin{equation}\label{eq:Momn}
% \mathbf{M}_{omn}
% =  \frac{m}{\sin\theta}\cos m\phi\, P_n^m(\cos\theta)\, z_n(\rho)\,\hat{\boldsymbol{e}}_\theta
%   - \sin m\phi\, \frac{d P_n^m(\cos\theta)}{d\theta}\, z_n(\rho)\,\hat{\boldsymbol{e}}_\phi ,
% \end{equation}

% \begin{equation}\label{eq:Nemn}
% \begin{split}
% \mathbf{N}_{emn} =
% & \frac{z_n(\rho)}{\rho}\cos m\phi\, n(n+1) P_n^m(\cos\theta)\,\hat{\mathbf{e}}_r \\
% & + \cos m\phi\, \frac{d P_n^m(\cos\theta)}{d\theta}\,
%     \frac{1}{\rho}\frac{d}{d\rho}\!\left[\rho z_n(\rho)\right]\hat{\mathbf{e}}_\theta \\
% & - m\sin m\phi\, \frac{P_n^m(\cos\theta)}{\sin\theta}\,
%     \frac{1}{\rho}\frac{d}{d\rho}\!\left[\rho z_n(\rho)\right]\hat{\mathbf{e}}_\phi ,
% \end{split}
% \end{equation}

% \begin{equation}\label{eq:Nomn}
% \begin{split}
% \mathbf{N}_{omn} =
% & \frac{z_n(\rho)}{\rho}\sin m\phi\, n(n+1) P_n^m(\cos\theta)\,\hat{\mathbf{e}}_r \\
% & + \sin m\phi\, \frac{d P_n^m(\cos\theta)}{d\theta}\,
%     \frac{1}{\rho}\frac{d}{d\rho}\!\left[\rho z_n(\rho)\right]\hat{\mathbf{e}}_\theta \\
% & + m\cos m\phi\, \frac{P_n^m(\cos\theta)}{\sin\theta}\,
%     \frac{1}{\rho}\frac{d}{d\rho}\!\left[\rho z_n(\rho)\right]\hat{\mathbf{e}}_\phi .
% \end{split}
% \end{equation}


Cualquier solución de las ecuaciones de campo puede expandirse como una serie infinita 
de las funciones (\ref{Generadores vectoriales}). En particular, el campo incidente puede 
escribirse como

\begin{equation}
\mathbf{E}_i = 
\sum_{m=0}^{\infty} \sum_{n=m}^{\infty}
\left(
B_{emn}\,\mathbf{M}_{emn}
+ B_{omn}\,\mathbf{M}_{omn}
+ A_{emn}\,\mathbf{N}_{emn}
+ A_{omn}\,\mathbf{N}_{omn}
\right).
\end{equation}



y los campos interno y dispersado admiten expansiones análogas con coeficientes 
distintos. Cada término satisface de forma automática la ecuación de onda en su 
dominio y, gracias a la ortogonalidad de los armónicos esféricos vectoriales, las 
condiciones de contorno en $r=a$ se reducen a un conjunto de ecuaciones 
algebraicas independientes. Esto posibilita obtener los coeficientes de Mie y, 
a partir de ellos, las propiedades ópticas de la partícula.


% Una frase que me gusta: Añadir más adelante:
% 2025-12-18. Mirar esto que está comentado. revisar simplemente.
% En problemas de dispersión por partículas esféricas resulta conveniente expresar 
% los campos electromagnéticos en una base adaptada a la simetría del sistema. 
% Los campos deben satisfacer la ecuación de onda vectorial
% \[
% \nabla^{2}\mathbf{E} + k^{2}\mathbf{E} = 0,
% \qquad 
% \nabla \cdot \mathbf{E} = 0,
% \]
% por lo que es natural buscar soluciones separables en coordenadas esféricas. 
% En este contexto, los armónicos esféricos vectoriales (VSHs) constituyen una 
% familia completa y ortogonal de funciones que resuelven la ecuación de onda y 
% cumplen la condición de transversidad.

% Los VSHs, habitualmente escritos como los modos magnéticos 
% $\mathbf{M}_{\ell m}$ y eléctricos $\mathbf{N}_{\ell m}$, permiten expresar 
% cualquier campo incidente, interno o dispersado como una expansión multipolar. 
% Cada par de índices $(\ell,m)$ define un modo con estructura angular fija, 
% mientras que la dependencia radial se describe mediante funciones de Bessel o 
% Hankel, según corresponda al interior o exterior de la partícula. 
% Esta formulación no solo hace el problema matemáticamente tratable, sino que 
% también revela explícitamente las resonancias y contribuciones multipolares 
% que dependen del tamaño, la permitividad y la frecuencia del sistema.
% Revisar este fragmento a ver si se puede incluir en algún sitio.

\subsubsection{Expansión de una onda plana en armónicos esféricos vectoriales}
\hspace{1em}Una vez definidos los armónicos esféricos vectoriales, es posible expandir
cualquier campo electromagnético en esta base. Ahora, se muestra cómo expresar una
onda plana en términos de los armónicos esféricos vectoriales, es decir,
se buscan los coeficientes $A_{emn}$, $A_{omn}$, $B_{emn}$ y $B_{omn}$
de la expansión.

Consideremos una onda plana monocromática, linealmente polarizada en $x$, que se 
propaga en la dirección $z$:
\begin{equation}
  \mathbf{E}_{i} = E_0 e^{ikz}\,\hat{\mathbf{e}}_x .
\end{equation}
donde $\hat{\mathbf{e}}_x = \sin\theta \cos \phi\hat{\mathbf{e}}_r + 
        \cos\theta \cos \phi\hat{\mathbf{e}}_\theta - \sin\phi\hat{\mathbf{e}}_\phi$.

\begin{equation}
  \mathbf{E}_{i} =
  \sum_{n=1}^{\infty}\sum_{m=-n}^{n}
  \left(
    B_{emn}\,\mathbf{M}_{emn}
  + B_{omn}\,\mathbf{M}_{omn}
  + A_{emn}\,\mathbf{N}_{emn}
  + A_{omn}\,\mathbf{N}_{omn}
  \right).
\end{equation}

% Como $\sin(m\phi)$ es ortogonal a $\cos(m'\phi)$ para todo $m$ y $m'$, las parejas 
% ($M_{emn}$,$M_omn$) y ($N_{omn}$,$N_{emn}$), ($M_{omn}$,$N_{omn}$) 
% y ($N_{emn}$,$N_{emn}$) son ortogonales entre sí.

Todo el set de funciones son ortogonales entre sí (pagina 90 del bohren), por lo que
los coeficientes de la expansión tienen la forma
\begin{equation}  
  B_{emn} =
  \frac{
    \int_{0}^{2\pi}\int_{0}^{\pi}
    \mathbf{E}_{i} \cdot \mathbf{M}_{emn} \sin\theta d\theta d\phi
  }{
    \int_{0}^{2\pi}\int_{0}^{\pi}
    |\mathbf{M}_{emn}|^2 \sin\theta d\theta d\phi
  },
\end{equation}

Para todo $m$ y $n$, se obtiene que $B_{emn} = A_{omn} = 0$ debido a la ortogonalidad
del seno y coseno y, por la misma razón, solo los términos con $m=1$ contribuyen a 
la expansión en $B_{omn}$ y $A_{emn}$. Además, la regularidad del campo en el origen exige que la dependencia radial 
de los armónicos generadores esté dada por las funciones de Bessel esféricas 
$j_n(kr)$. Introduciendo esta condición, la expansión del campo incidente se 
simplifica y adopta la forma


\begin{equation}
  \mathbf{E}_{i} =
  \sum_{n=1}^{\infty}
    \left(
    B_{o1n}\mathbf{M}_{o1n}^{(1)}
    + A_{e1n}\mathbf{N}_{e1n}^{(1)}
  \right),
\end{equation}

donde el superíndice $(1)$ indica que la dependencia radial viene dada por las 
funciones de Bessel esféricas de primera especie/tipo (kind?). Finalmente, los 
coeficientes de la expansión se obtienen evaluando las integrales 
anteriores (Bohren pag 92), logrando para el campo eléctrico y magnético incidente:


\begin{equation}
  \mathbf{E}_{i} =
  E_0 \sum_{n=1}^{\infty}
  i^n \frac{2n+1}{n(n+1)}
  \left(
    \mathbf{M}_{o1n}^{(1)}
    - i\,\mathbf{N}_{e1n}^{(1)}
  \right),
\end{equation}

\begin{equation}
  \mathbf{H}_{i} =
  -\frac{k}{\omega \mu}
  E_0 \sum_{n=1}^{\infty}
  i^n \frac{2n+1}{n(n+1)}
  \left(
    \mathbf{M}_{e1n}^{(1)}
    + i\,\mathbf{N}_{o1n}^{(1)}
  \right).
\end{equation}



Esta expresión proporciona la expansión completa de una onda plana en términos 
de armónicos esféricos vectoriales y constituye el punto de partida para la 
aplicación de las condiciones de contorno en la superficie de la esfera. A partir 
de esta expansión, los campos interno y dispersado pueden construirse de manera 
análoga, permitiendo la determinación sistemática de los coeficientes de Mie.



\subsubsection{Campos internos y dispersados}
\hspace{1em}En el apartado anterior se obtuvo la expansión del campo incidente en armónicos
esféricos vectoriales. El siguiente paso consiste en expresar los campos interno y
dispersado. En la frontera de la esfera, imponemos las condiciones de contorno

\begin{equation}
  \left( \mathbf{E}_{i} + \mathbf{E}_{sca} - \mathbf{E}_{int} \right) \times \hat{\mathbf{e}}_r
=
\left( \mathbf{H}_{i} + \mathbf{H}_{sca} - \mathbf{H}_{int} \right) \times \hat{\mathbf{e}}_r
= 0.
\end{equation}

Las condiciones de contorno, la ortogonalidad de los armónicos esféricos vectoriales y la
forma de la expansión del campo incidente implican que los campos interno y dispersado
tienen la forma siguiente:

\begin{equation}
  \mathbf{E}_{int} =
  \sum_{n=1}^{\infty}
    E_{n}\left(
    c_{n}\mathbf{M}_{o1n}^{(1)}
    - i d_{n}\mathbf{N}_{e1n}^{(1)}
  \right)
\end{equation}
\begin{equation}
  \mathbf{H}_{int} =
  -\frac{k_{1}}{\omega \mu_{1}}
  \sum_{n=1}^{\infty}
    E_{n}\left(
    d_{n}\mathbf{M}_{e1n}^{(1)}
    + i c_{n}\mathbf{N}_{o1n}^{(1)}
  \right)
\end{equation}

donde $E_{n} = i^{n} E_{0} (2n+1)/[n(n+1)]$, $k_{1}$ y $\mu_{1}$ son el
número de onda y la permeabilidad del interior de la esfera, y
$c_{n}$ y $d_{n}$ son los coeficientes de campo interno.

\begin{equation}
  \mathbf{E}_{sca} =
  \sum_{n=1}^{\infty}
    E_{n}\left(
    i a_{n}\mathbf{N}_{e1n}^{(3)}
    - b_{n}\mathbf{M}_{o1n}^{(3)}
  \right)
\end{equation}
\begin{equation}
  \mathbf{H}_{sca} =
  \frac{k}{\omega \mu}
  \sum_{n=1}^{\infty}
    E_{n}\left(
    i b_{n}\mathbf{N}_{o1n}^{(3)}
    + a_{n}\mathbf{M}_{e1n}^{(3)}
  \right)
\end{equation}
donde los superíndices $(3)$ indican que la dependencia radial está dada por las 
funciones de Hankel esféricas de primera especie/tipo, y $a_{n}$ y $b_{n}$ son los
coeficientes de scattering. 

\subsubsection{Coeficientes de scattering}
El siguiente paso nos lleva a conseguir las expresiones explícitas de los coeficientes
$a_{n}$, $b_{n}$, $c_{n}$ y $d_{n}$. Estos se obtienen imponiendo las condiciones de
contorno en la superficie de la esfera, que exigen la continuidad de las componentes
tangenciales de los campos eléctrico y magnético en $r=a$:

\[
\begin{aligned}
E_{i\theta} + E_{s\theta} &= E_{1\theta}, \\
E_{i\phi} + E_{s\phi} &= E_{1\phi}, \qquad r = a, \\
H_{i\theta} + H_{s\theta} &= H_{1\theta}, \\
H_{i\phi} + H_{s\phi} &= H_{1\phi}.
\end{aligned}
\]

Al sustituir las expansiones de los campos en estas ecuaciones y utilizar la
ortogonalidad de los armónicos esféricos vectoriales, se obtiene un sistema de ecuaciones
lineales para cada valor de $n$.

\begin{equation}
x = ka = \frac{2\pi N a}{\lambda}, 
\qquad
m = \frac{k_1}{k} = \frac{N_1}{N}.
\end{equation}


Aquí, $N_1$ y $N$ son los índices de refracción de la partícula y del medio,
respectivamente. Las cuatro ecuaciones lineales simultáneas se resuelven
para los coeficientes del campo dentro de la partícula:


\begin{align}
c_n &=
\frac{
\mu_1\, j_n(x)\,\left[x\, h_n^{(1)}(x)\right]^{\prime}
-
\mu_1\, h_n^{(1)}(x)\,\left[x\, j_n(x)\right]^{\prime}
}{
\mu\, j_n(mx)\,\left[x\, h_n^{(1)}(x)\right]^{\prime}
-
\mu_1\, h_n^{(1)}(x)\,\left[mx\, j_n(mx)\right]^{\prime}
},
\\[1em]
d_n &=
\frac{
\mu_1\, m\, j_n(x)\,\left[x\, h_n^{(1)}(x)\right]^{\prime}
-
\mu_1\, m\, h_n^{(1)}(x)\,\left[x\, j_n(x)\right]^{\prime}
}{
\mu\, m^2\, j_n(mx)\,\left[x\, h_n^{(1)}(x)\right]^{\prime}
-
\mu_1\, h_n^{(1)}(x)\,\left[mx\, j_n(mx)\right]^{\prime}
}.
\end{align}

Los coeficientes de dispersión vienen dados por

\begin{align}
a_n &=
\frac{
\mu\, m^2\, j_n(mx)\,\left[x\, j_n(x)\right]^{\prime}
-
\mu_1\, j_n(x)\,\left[mx\, j_n(mx)\right]^{\prime}
}{
\mu\, m^2\, j_n(mx)\,\left[x\, h_n^{(1)}(x)\right]^{\prime}
-
\mu_1\, h_n^{(1)}(x)\,\left[mx\, j_n(mx)\right]^{\prime}
},
\\[1em]
b_n &=
\frac{
\mu_1\, j_n(mx)\,\left[x\, j_n(x)\right]^{\prime}
-
\mu\, j_n(x)\,\left[mx\, j_n(mx)\right]^{\prime}
}{
\mu_1\, j_n(mx)\,\left[x\, h_n^{(1)}(x)\right]^{\prime}
-
\mu\, h_n^{(1)}(x)\,\left[mx\, j_n(mx)\right]^{\prime}
}.
\end{align}


%====================================================================%
%====================================================================%

\subsection{Concepto físico de dispersión}

\hspace{1em}Desde un punto de vista físico, la dispersión puede interpretarse 
como la rerradiación de energía electromagnética por parte de las cargas ligadas 
y corrientes de polarización inducidas en la partícula por el campo incidente. 
La onda incidente excita momentos multipolares en el interior de la partícula, 
que a su vez emiten radiación en todas las direcciones. En medios disipativos, 
parte de la energía electromagnética se transforma en energía interna del 
material, dando lugar a la absorción. El balance energético de la interacción 
se expresa mediante la relación

\begin{equation}
  C_{\text{ext}} = C_{\text{sca}} + C_{\text{abs}},
\end{equation}

donde $C_{\text{ext}}$ es la sección eficaz de extinción, $C_{\text{sca}}$ la 
de dispersión y $C_{\text{abs}}$ la de absorción. Estas magnitudes cuantifican 
el área efectiva que presenta la partícula frente al haz incidente.


%====================================================================%
\subsection{Secciones eficaces}
%====================================================================%

\hspace{1em}Las secciones eficaces se obtienen aplicando el teorema óptico y 
el teorema de Poynting al campo total. El flujo de energía dispersada y 
extinguida por la partícula se calcula integrando el vector de Poynting sobre 
una superficie esférica que la rodea, y proyectando sobre las expansiones 
modales de los campos. Gracias a la ortogonalidad de los armónicos esféricos 
vectoriales, estas integrales se reducen a sumas sobre los coeficientes de 
Mie, dando lugar a las expresiones

\begin{equation}
  C_{\text{sca}} = \frac{2\pi}{k^2}
  \sum_{n=1}^{\infty} (2n+1)
  \left( |a_n|^2 + |b_n|^2 \right),
\end{equation}

\begin{equation}
  C_{\text{ext}} = \frac{2\pi}{k^2}
  \sum_{n=1}^{\infty} (2n+1)
  \,\mathrm{Re}(a_n + b_n),
\end{equation}

\begin{equation}
  C_{\text{abs}} = C_{\text{ext}} - C_{\text{sca}}.
\end{equation}

La sección eficaz de extinción sigue el \textit{teorema óptico}: está 
determinada únicamente por la amplitud de dispersión en la dirección hacia 
adelante ($\theta = 0$), lo que refleja la interferencia destructiva entre 
el campo incidente y el dispersado que produce la sombra de la partícula.

Es habitual introducir las eficiencias adimensionales, definidas como las 
secciones eficaces normalizadas por el área geométrica de la partícula 
$G = \pi a^2$:

\begin{equation}
  Q_{\text{sca}} = \frac{C_{\text{sca}}}{\pi a^2}, \qquad
  Q_{\text{ext}} = \frac{C_{\text{ext}}}{\pi a^2}, \qquad
  Q_{\text{abs}} = \frac{C_{\text{abs}}}{\pi a^2}.
\end{equation}

En el límite de partícula grande ($x \gg 1$), $Q_{\text{ext}} \to 2$, 
resultado conocido como la \textit{paradoja de extinción}: la partícula 
extingue el doble de la energía que intercepta geométricamente, debido 
a la difracción en su borde.


%====================================================================%
\subsection{Elementos de la matriz de amplitud de dispersión}
%====================================================================%

\hspace{1em}A partir de los coeficientes de Mie $a_n$ y $b_n$ se obtienen 
los elementos de la matriz de amplitud de dispersión introducida en 
(\ref{eq:amplitudematrix}). Para ello se evalúa el campo dispersado 
(ref{}) en la zona de campo lejano y se proyectan sus componentes de 
polarización. El resultado son las funciones

\begin{equation}
\label{eq:S1S2}
  S_1(\theta) = \sum_{n=1}^{\infty} \frac{2n+1}{n(n+1)}
  \left[ a_n\, \pi_n(\cos\theta) + b_n\, \tau_n(\cos\theta) \right],
\end{equation}

\begin{equation}
  S_2(\theta) = \sum_{n=1}^{\infty} \frac{2n+1}{n(n+1)}
  \left[ b_n\, \pi_n(\cos\theta) + a_n\, \tau_n(\cos\theta) \right],
\end{equation}

donde las \textit{funciones angulares de Mie} se definen como

\begin{equation}
  \pi_n(\cos\theta) = \frac{P_n^1(\cos\theta)}{\sin\theta},
  \qquad
  \tau_n(\cos\theta) = \frac{d\,P_n^1(\cos\theta)}{d\theta},
\end{equation}

siendo $P_n^1$ el polinomio de Legendre asociado de orden $m=1$. Estas 
funciones satisfacen las relaciones de recurrencia

\begin{equation}
  \pi_n = \frac{2n-1}{n-1}\cos\theta\,\pi_{n-1} - \frac{n}{n-1}\pi_{n-2},
  \qquad
  \tau_n = n\cos\theta\,\pi_n - (n+1)\pi_{n-1},
\end{equation}

con valores iniciales $\pi_0 = 0$, $\pi_1 = 1$, que permiten su evaluación 
numérica eficiente.

Para una esfera homogénea e isotrópica, $S_3 = S_4 = 0$ por simetría 
azimutal, de modo que la matriz de amplitud es diagonal y los modos de 
polarización paralelo y perpendicular se dispersan de forma independiente. 
En el límite de dispersión hacia adelante ($\theta = 0$) se recupera 
el teorema óptico en su forma explícita:

\begin{equation}
  C_{\text{ext}} = \frac{4\pi}{k^2}\,\mathrm{Re}\left[S_1(0)\right]
                 = \frac{4\pi}{k^2}\,\mathrm{Re}\left[S_2(0)\right],
\end{equation}

siendo $S_1(0) = S_2(0)$ por la isotropía de la esfera.


%====================================================================%
\subsection{Matriz de Mueller}
%====================================================================%

\hspace{1em}Para relacionar las \textit{intensidades} y el estado de 
polarización del campo dispersado con los del campo incidente, se introduce 
la \textit{matriz de Mueller} $\mathbf{M}(\theta)$, que actúa sobre el 
vector de Stokes $(I, Q, U, V)^T$:

\begin{equation}
  \begin{pmatrix} I^{(s)} \\ Q^{(s)} \\ U^{(s)} \\ V^{(s)} \end{pmatrix}
  =
  \frac{1}{k^2 r^2}
  \begin{pmatrix}
    S_{11} & S_{12} & 0      & 0      \\
    S_{12} & S_{11} & 0      & 0      \\
    0      & 0      & S_{33} & S_{34} \\
    0      & 0      &-S_{34} & S_{33}
  \end{pmatrix}
  \begin{pmatrix} I^{(i)} \\ Q^{(i)} \\ U^{(i)} \\ V^{(i)} \end{pmatrix}.
\end{equation}

Los elementos de la matriz de Mueller se construyen como productos bilineales 
de $S_1$ y $S_2$:

\begin{equation}
  S_{11} = \tfrac{1}{2}\left(|S_2|^2 + |S_1|^2\right), \qquad
  S_{12} = \tfrac{1}{2}\left(|S_2|^2 - |S_1|^2\right),
\end{equation}

\begin{equation}
  S_{33} = \tfrac{1}{2}\left(S_2 S_1^* + S_2^* S_1\right), \qquad
  S_{34} = \tfrac{i}{2}\left(S_1 S_2^* - S_2 S_1^*\right).
\end{equation}

La estructura en bloque de la matriz refleja la simetría de la esfera: 
los bloques diagonales $2\times 2$ de $U$ y $V$ están desacoplados de $I$ 
y $Q$. El elemento $S_{11}(\theta)$ es la \textit{función de fase}, que 
describe la distribución angular de la energía dispersada y es directamente 
medible. La relación $S_{11}^2 \geq S_{12}^2 + S_{33}^2 + S_{34}^2$ 
garantiza la positividad física de la matriz.


%====================================================================%
\subsection{Función de fase y normalización}
%====================================================================%

\hspace{1em}El elemento $S_{11}(\theta)$ de la matriz de Mueller describe 
la distribución angular de la energía dispersada. Se define 
la \textit{función de fase} $p(\theta)$ como

\begin{equation}
  p(\theta) = \frac{4\pi}{k^2\, C_{\text{sca}}}\, S_{11}(\theta),
\end{equation}


y es la distribución angular normalizada 
de la potencia dispersada por unidad de ángulo sólido. Satisface la 
condición de normalización

\begin{equation}
  \frac{1}{4\pi} \int_{4\pi} p(\theta)\, d\Omega 
  = \frac{1}{2} \int_0^{\pi} p(\theta)\sin\theta\, d\theta = 1,
\end{equation}

de modo que $p(\theta)\,d\Omega / 4\pi$ representa la fracción de la 
potencia total dispersada que se redistribuye en el elemento de ángulo 
sólido $d\Omega$ en torno a la dirección $\theta$. Esta normalización 
facilita la comparación entre partículas de distinto tamaño y constituye 
la forma estándar de caracterizar la anisotropía angular de la dispersión.


%====================================================================%
\subsection{Parámetro de asimetría}
%====================================================================%

\hspace{1em}La función de fase contiene toda la información angular de la 
dispersión, pero en muchas aplicaciones resulta útil resumirla en un único 
parámetro escalar. El \textit{parámetro de asimetría} $g$ se define como 
el valor medio del coseno del ángulo de dispersión ponderado por la función 
de fase:

\begin{equation}
  g = \langle \cos\theta \rangle 
    = \frac{1}{2}\int_0^{\pi} p(\theta)\cos\theta\sin\theta\, d\theta
    = \frac{\int_0^\pi S_{11}(\theta)\cos\theta\sin\theta\,d\theta}
           {\int_0^\pi S_{11}(\theta)\sin\theta\,d\theta}.
\end{equation}

El parámetro $g$ toma valores en el intervalo $[-1, 1]$. Un valor $g > 0$ 
indica que la partícula dispersa preferentemente en la dirección hacia 
adelante, mientras que $g < 0$ corresponde a dispersión preferente hacia 
atrás. El caso $g = 0$ corresponde a dispersión isótropa. En términos 
de los coeficientes de Mie, $g$ puede expresarse analíticamente como

\begin{equation}
  g = \frac{4}{Q_{\text{sca}}\,x^2}
  \left[
    \sum_{n=1}^{\infty}
    \frac{n(n+2)}{n+1}\,\mathrm{Re}(a_n a_{n+1}^* + b_n b_{n+1}^*)
    +
    \sum_{n=1}^{\infty}
    \frac{2n+1}{n(n+1)}\,\mathrm{Re}(a_n b_n^*)
  \right].
\end{equation}

Para partículas pequeñas en el régimen de Rayleigh ($x \ll 1$), la 
dispersión tiende a ser casi isótropa y $g \approx 0$. A medida que 
el parámetro de tamaño aumenta, la dispersión se vuelve cada vez más 
anisótropa hacia adelante y $g \to 1$.


%====================================================================%
\subsection{Albedo de dispersión simple}
%====================================================================%

\hspace{1em}De la energía total extinguida por la partícula, una fracción 
es redistribuida espacialmente mediante dispersión y el resto es absorbida 
y convertida en energía interna. El \textit{albedo de dispersión simple} 
$\omega_0$ cuantifica esta fracción:

\begin{equation}
  \omega_0 = \frac{C_{\text{sca}}}{C_{\text{ext}}} 
           = \frac{Q_{\text{sca}}}{Q_{\text{ext}}}.
\end{equation}

Por definición, $\omega_0 \in [0,1]$. El límite $\omega_0 = 1$ corresponde 
a una partícula puramente dispersora (sin absorción, $C_{\text{abs}} = 0$), 
mientras que $\omega_0 = 0$ corresponde a una partícula puramente absorbente. 
En términos del índice de refracción, una parte imaginaria nula del índice 
($k = 0$, medio transparente) conduce a $\omega_0 = 1$, mientras que una 
parte imaginaria elevada implica $\omega_0 \ll 1$.

Los tres parámetros $\{Q_{\text{ext}},\, \omega_0,\, g\}$ constituyen la 
descripción óptica mínima de una partícula individual y son precisamente 
las magnitudes de entrada que requiere la ecuación de transporte radiativo 
para describir un medio dispersor.


%====================================================================%
\subsection{Regímenes de dispersión}
%====================================================================%

\hspace{1em}El comportamiento óptico de una partícula esférica depende 
fundamentalmente del \textit{parámetro de tamaño} $x = 2\pi a / \lambda$, 
que compara la circunferencia de la partícula con la longitud de onda 
incidente. En función del valor de $x$ se distinguen tres regímenes 
cualitativamente distintos.

\subsubsection{Régimen de Rayleigh ($x \ll 1$)}

Cuando la partícula es mucho más pequeña que la longitud de onda, el campo 
incidente varía poco sobre su volumen y la respuesta queda dominada por el 
término dipolar ($n=1$). En este límite, los coeficientes de Mie se 
aproximan por

\begin{equation}
  a_1 \approx -\frac{2i x^3}{3}\frac{m^2-1}{m^2+2}, 
  \qquad
  b_1 \approx 0,
\end{equation}

y los de orden superior son despreciables. Las secciones eficaces escalan como

\begin{equation}
  C_{\text{sca}} \propto x^4 \propto \lambda^{-4},
  \qquad
  C_{\text{abs}} \propto x \propto \lambda^{-1},
\end{equation}

lo que explica el color azul del cielo: la dispersión preferente de las 
longitudes de onda cortas por las moléculas de aire (partículas en régimen 
de Rayleigh) da lugar a la dispersión $\propto \lambda^{-4}$.
La dispersión es casi isótropa ($g \approx 0$) con un débil lóbulo 
hacia adelante y hacia atrás simétricos.

\subsubsection{Régimen de Mie ($x \sim 1$)}

Cuando el tamaño de la partícula es comparable a la longitud de onda, 
múltiples órdenes multipolares contribuyen de forma significativa. 
Las secciones eficaces presentan oscilaciones resonantes en función de $x$, 
conocidas como \textit{resonancias de Mie}, asociadas a los modos 
electromagnéticos del interior de la esfera. La dispersión se vuelve 
marcadamente anisótropa hacia adelante ($g \to 1$) y la función de fase 
desarrolla lóbulos bien definidos. Este es el régimen de mayor riqueza 
fenomenológica y el que requiere el tratamiento completo de la teoría de Mie.

\subsubsection{Régimen de óptica geométrica ($x \gg 1$)}

Cuando la partícula es mucho mayor que la longitud de onda, los efectos 
de difracción se concentran en un ángulo muy estrecho alrededor de la 
dirección de propagación y la interacción puede describirse mediante 
rayos geométricos. En este límite, la eficiencia de extinción converge a

\begin{equation}
  Q_{\text{ext}} \xrightarrow{x \to \infty} 2,
\end{equation}

resultado conocido como la \textit{paradoja de extinción}: la partícula 
extingue el doble de la energía que intercepta geométricamente. El factor 
adicional proviene de la difracción de Fraunhofer en el borde de la 
partícula, que desvía energía fuera del haz hacia adelante con la misma 
magnitud que la energía interceptada.


%====================================================================%
\subsection{Parámetros de Stokes y grado de polarización}
%====================================================================%

\hspace{1em}El estado de polarización de una onda electromagnética queda 
completamente caracterizado por los \textit{parámetros de Stokes}, definidos 
en términos de las componentes del campo eléctrico como

\begin{equation}
\begin{aligned}
  I &= \langle E_\parallel E_\parallel^* \rangle 
     + \langle E_\perp    E_\perp^*    \rangle, \\
  Q &= \langle E_\parallel E_\parallel^* \rangle 
     - \langle E_\perp    E_\perp^*    \rangle, \\
  U &= \langle E_\parallel E_\perp^*    \rangle 
     + \langle E_\perp    E_\parallel^* \rangle, \\
  V &= i\left(
         \langle E_\perp    E_\parallel^* \rangle 
       - \langle E_\parallel E_\perp^*   \rangle
       \right),
\end{aligned}
\end{equation}

donde los corchetes indican promedio temporal. A diferencia de los campos, 
los parámetros de Stokes son magnitudes reales y directamente medibles 
mediante combinaciones de polarizadores y retardadores. Satisfacen la 
relación

\begin{equation}
  I^2 \geq Q^2 + U^2 + V^2,
\end{equation}

donde la igualdad se cumple para luz completamente polarizada. El 
\textit{grado de polarización total} se define como

\begin{equation}
  P = \frac{\sqrt{Q^2 + U^2 + V^2}}{I},
\end{equation}

y toma valores en $[0,1]$. Para luz natural (no polarizada), $Q = U = V = 0$ 
y $P = 0$. Para cuantificar únicamente la componente lineal de la 
polarización se introduce el \textit{grado de polarización lineal}

\begin{equation}
  P_L = \frac{\sqrt{Q^2 + U^2}}{I}.
\end{equation}

En el caso particular de luz incidente no polarizada que interactúa con 
una esfera homogénea, $U = V = 0$ en el campo dispersado y el grado de 
polarización lineal toma la forma

\begin{equation}
  P_L(\theta) = -\frac{S_{12}(\theta)}{S_{11}(\theta)},
\end{equation}

que es directamente medible experimentalmente y proporciona información 
sobre la composición y el tamaño de las partículas independientemente de 
su concentración.


\cite{Bohren}\cite{Mishchenko}\cite{Bohren,Mishchenko} (prueba, funcionan las citas)
\cite{Bohren}
